\section{Method}\label{sec:method}

\subsection{Crisp RFM-FCA}\label{sec:crisp}

Let $G$ be a set of customers. For each customer, recency $R$ (days since the last purchase), frequency $F$ (number of transaction lines) and monetary value $M$ (total spend) are computed at a cutoff date. Each dimension $d \in \{R,F,M\}$ is mapped to a mid-rank percentile $u_d(g) \in (0,1]$ among the fitted customers, oriented so that larger is better (recency is reversed). The crisp quintile score is $q_d(g) = \lceil 5\,u_d(g) \rceil$, and the crisp context is the binary relation $I \subseteq G \times M$ over the fifteen attributes $M = \{R1,\dots,R5,F1,\dots,F5,M1,\dots,M5\}$ with $(g, dk) \in I$ iff $q_d(g) = k$. Formal concepts are pairs $(A,B)$ with $A' = B$ and $B' = A$ under the usual derivation operators \citep{ganter1999}; a concept's \emph{support} is $|A|/|G|$, and \citet{rungruang2024} keep concepts with support at least $0.04$ as segments.

\subsection{Fuzzy RFM context}\label{sec:fuzzy-context}

\begin{definition}[Linguistic terms]
For $k=1,\dots,5$ let $c_k = (2k-1)/10$ be the quintile midpoints and $h = 0.2$ the half-width. The membership of rank $u$ in term $k$ is
\begin{equation}
\mu_k(u) = \max\!\left(0,\; 1 - \frac{|u - c_k|}{h}\right),
\end{equation}
with saturating shoulders $\mu_1(u) = 1$ for $u \le c_1$ and $\mu_5(u) = 1$ for $u \ge c_5$.
\end{definition}

The five terms form a Ruspini partition \citep{ruspini1969}: $\sum_k \mu_k(u) = 1$ and at most two adjacent terms are positive. A customer in the middle of a quintile band is fully in that term; a customer near a cut point is partly in the two neighbouring terms.

Memberships are rounded to the chain $L = \{0, \tfrac12, 1\}$ (``not'', ``partly'', ``fully''), which keeps the lattice finite and the segment descriptions readable. The fuzzy context is $\tilde I(g, dk) = \operatorname{round}_L \mu_k(u_d(g))$.

\subsection{Fuzzy concepts}\label{sec:fuzzy-concepts}

We use the fuzzy concept lattice of \citet{belohlavek2002,belohlavek2004} over the chain $L$ with a residuated implication $\to$; our primary choice is the Łukasiewicz residuum $a \to b = \min(1, 1-a+b)$, with the Gödel residuum ($a \to b = 1$ if $a \le b$, else $b$) as a check \citep{hajek1998}. For a fuzzy set of customers $A$ and of attributes $B$,
\begin{equation}
A^{\uparrow}(m) = \bigwedge_{g \in G} \bigl(A(g) \to \tilde I(g,m)\bigr), \qquad
B^{\downarrow}(g) = \bigwedge_{m \in M} \bigl(B(m) \to \tilde I(g,m)\bigr).
\end{equation}
A fuzzy concept is a pair $(A,B)$ with $A^{\uparrow} = B$ and $B^{\downarrow} = A$; $B^{\downarrow}(g)$ is customer $g$'s degree of membership in the concept. Since $A^{\uparrow}$ is the meet of the intents $\{a/g\}^{\uparrow} = a \to \tilde I(g,\cdot)$, the set of intents is the meet-closure of these generators together with the top element; we compute it by iterated pointwise minima over the distinct rows of $\tilde I$.

\begin{definition}[Core support]
The \emph{core extent} of a concept with intent $B$ is the crisp set $\{g : B^{\downarrow}(g) = 1\}$ and its \emph{core support} is $\mathrm{cs}(B) = |\{g : B^{\downarrow}(g) = 1\}| / |G|$. Concepts with $\mathrm{cs}(B) \ge 0.04$ (excluding the top concept) are the segments of \emph{fuzzy RFM-FCA}; customers are represented by their degrees $B^{\downarrow}(g)$ in these segments.
\end{definition}

Rules between segments are read along the covers of the order restricted to the segments, with confidence $\mathrm{cs}(B_2)/\mathrm{cs}(B_1)$ for $B_1 < B_2$, in the spirit of the Luxenburger basis \citep{luxenburger1991}.

\subsection{Properties}\label{sec:properties}

\begin{proposition}[Crisp RFM-FCA is a special case]\label{prop:special}
(i) If $u$ lies in the interior of quintile band $q$, then $\mu_q(u) > \tfrac12 > \mu_k(u)$ for all $k \ne q$; hence the 0.5-cut of the fuzzy context is the crisp quintile context. (ii) Over the two-element chain $\{0,1\}$ every residuated implication is classical implication, so the fuzzy concept lattice of a binary context equals its classical concept lattice, degrees are extent indicators, and core support equals support.
\end{proposition}
\begin{proof}
(i) Band $q$ is $((q-1)/5, q/5]$ with midpoint $c_q$; an interior point is less than $0.1$ from $c_q$ and more than $0.1$ from every other midpoint, so $\mu_q(u) > 1 - 0.1/0.2 = \tfrac12$ and $\mu_k(u) < \tfrac12$. (ii) On $\{0,1\}$ residuation forces $a \to b = 1$ iff $a \le b$, i.e.\ $0\to 0 = 0\to1 = 1\to1 = 1$ and $1 \to 0 = 0$; the operators $\uparrow,\downarrow$ then coincide with the derivation operators.
\end{proof}

\begin{proposition}[Core support]\label{prop:core}
For any residuated implication on a chain: (i) $B^{\downarrow}(g) = 1$ iff $B(m) \le \tilde I(g,m)$ for all $m$, so the core extent does not depend on the residuum; (ii) core support is anti-monotone: $B_1 \le B_2$ implies $\mathrm{cs}(B_2) \le \mathrm{cs}(B_1)$, so frequent intents form a down-set and can be mined level-wise; (iii) on a binary context core support equals ordinary support.
\end{proposition}
\begin{proof}
(i) In a residuated lattice $a \to b = 1$ iff $a \le b$, and a meet over a chain equals 1 iff every term does. (ii) If $B_1 \le B_2$ and $B_2(m) \le \tilde I(g,m)$ for all $m$, then $B_1(m) \le \tilde I(g,m)$. (iii) Follows from (i) and Proposition~\ref{prop:special}(ii).
\end{proof}

\begin{proposition}[Łukasiewicz floor]\label{prop:floor}
Under the Łukasiewicz residuum, $B^{\downarrow}(g) \ge 1 - \max_m B(m)$ for every customer. Consequently the sigma-count support $\sum_g B^{\downarrow}(g)/|G|$ of every intent whose degrees are all at most $\tfrac12$ is at least $\tfrac12$, and every such concept passes any support threshold $s \le \tfrac12$.
\end{proposition}
\begin{proof}
$a \to b = \min(1, 1-a+b) \ge 1-a$ for $b \ge 0$, and the meet over $m$ is at least $\min_m (1 - B(m)) = 1 - \max_m B(m)$.
\end{proof}

Proposition~\ref{prop:floor} explains the failure of the naive extension: on Online Retail II, 2,649 of 3,841 Łukasiewicz concepts reach a sigma-count support of $0.04$, dominated by intents such as $\{\tfrac12 F1, \tfrac12 F2, \tfrac12 M1, \tfrac12 M2\}$ with support around $0.52$. Core support is immune to the floor, and by Proposition~\ref{prop:core} it is the natural generalisation of the crisp threshold.

\subsection{Summary view}\label{sec:summary}

The full set of fuzzy segments is about four times larger than the crisp one (Section~\ref{sec:res-interp}). For communication we select $K$ segments. A customer's \emph{explained mass} given a set $S$ of segments is the largest intent size $w_c = \sum_m B_c(m)$ (a full term counts 1, a partial term $\tfrac12$) among segments in $S$ whose core contains the customer, or 0. The objective is
\begin{equation}
f(S) = \sum_{g \in G} \max_{c \in S} w_c \,\mathbf{1}\bigl[B_c^{\downarrow}(g) = 1\bigr], \qquad f(\emptyset) = 0,
\end{equation}
maximised greedily under $|S| = K$, with ties broken by larger core support. The selection uses no outcome variable.

\begin{proposition}[Greedy guarantee]\label{prop:greedy}
$f$ is non-negative, monotone and submodular; the greedy selection $S_K$ satisfies $f(S_K) \ge (1 - 1/e) \max_{|S| = K} f(S)$.
\end{proposition}
\begin{proof}
For each $g$, $S \mapsto \max_{c \in S} a_{gc}$ with $a_{gc} \ge 0$ is monotone and submodular (a facility-location term); sums of such functions inherit both properties. The bound is the classical result of \citet{nemhauser1978}.
\end{proof}

\subsection{Tie-aware memberships}\label{sec:ties}

When many customers share a value (e.g.\ the 52--65\,\% of CDNOW customers with a single purchase), the mid-rank places the whole block at one point and the fuzzy context cannot express that its members' quintile is ambiguous. A block of $n_t > 1$ tied customers occupies the rank interval $[l, r]$ of width $n_t/|G|$ centred on the mid-rank. The tie-aware membership is
\begin{equation}
\bar\mu_k = \frac{1}{r - l} \int_l^r \mu_k(u)\,\mathrm{d}u ,
\end{equation}
with $\bar\mu_k = \mu_k(u)$ for untied customers.

\begin{proposition}[Tie-aware memberships]\label{prop:ties}
$\bar\mu$ is the expected membership when the customer's position within the tied block is uniform; it is again a partition ($\sum_k \bar\mu_k = 1$), and it coincides with $\mu$ for untied values.
\end{proposition}
\begin{proof}
Linearity of the integral and $\sum_k \mu_k \equiv 1$.
\end{proof}

All five propositions are checked numerically in the accompanying test suite.
